O ensino remoto é uma modalidade que pressupõe o distanciamento geográfico entre professores e alunos, adotada temporariamente nos diferentes níveis de ensino ao redor do mundo para evitar a interrupção das atividades escolares \cite{beharArtigoEnsinoRemoto2020}. A categoria é semelhante à Educação a Distância (EaD), mas sua origem está relacionada ao contexto da pandemia de COVID-19, configurando-se como uma solução emergencial para substituir o ensino presencial. No Brasil, essa modalidade foi aplicada em duas abordagens principais: a primeira consistiu em blocos de atividades, tarefas e estudos dirigidos, com acompanhamento por meio de mensagens; a segunda utilizou plataformas de webconferência para a realização de aulas síncronas \cite{santosEducacaoDistanciaEnsino2022}.

\cite{beharArtigoEnsinoRemoto2020} caracteriza o ensino remoto como a transposição da abordagem presencial para os meios digitais, ocorrendo de forma síncrona, por meio de webconferências, e com atividades de fixação assíncronas. A autora complementa que “a presença física do professor e do aluno na sala de aula é substituída por uma presença digital numa aula online, o que se chama de ‘presença social’, forma pela qual a tecnologia projeta essa presença”.

\cite{suhrReflexoesPapelAula2010} levanta a problemática da reprodução do modelo presencial no formato remoto sem adaptação metodológica. Esse movimento pode comprometer o potencial da modalidade, uma vez que essa abordagem permite explorar a diversidade inerente aos grupos sociais e produzir experiências educativas mais efetivas \cite{belloniEnsaioSobreEducacao2002}. A aplicação de metodologias centradas na exposição, marcadas pela comunicação unilateral e pelo protagonismo do professor, pode, assim, representar uma limitação conceitual no contexto do ensino remoto.

Muller (2009, \textit{apud} \cite{vasconcelosEstrategiasEnsinoAplicaveis2013}) ressalta que “tecnologias de informação e comunicação, em convergência com as telecomunicações e redes, mais especificamente a \textit{internet}, podem se tornar um meio de excelência para promover educação continuada e a distância”. No entanto, os desafios não podem ser ignorados: taxas elevadas de evasão, dispersão da atenção, dificuldades de acompanhamento individualizado e desigualdade no acesso às tecnologias configuram obstáculos recorrentes no contexto brasileiro.

Portanto, é essencial avaliar quais técnicas de ensino conseguem potencializar os recursos tecnológicos disponíveis, de modo a garantir não apenas a ampliação do acesso, mas também a qualidade e a permanência dos estudantes na modalidade a distância.